%!TEX program = xelatex 
\documentclass[a4paper,12pt]{article}
\usepackage{fontspec}\defaultfontfeatures{Ligatures=TeX}
\usepackage[a4paper,vmargin={3cm,3.5cm},hmargin={3cm,3cm}]{geometry}
%-----------------------------------------------------------------------------%
\usepackage{settings}
%-----------------------------------------------------------------------------%
%%% Title %%%
\title{G6411 Recitation Notes: Generalized Least Squares}
\author{Jesse Chieh Chen\\\href{mailto:cc5229@columbia.edu}{cc5229@columbia.edu}}
\date{\today}
%-----------------------------------------------------------------------------%

\begin{document}

\maketitle

\noindent
Heteroskedasticity and serial correlation need not make \gls{ols} inconsistent,
but they can make it inefficient.
Robust standard errors address inference without changing the coefficient estimates.
The classical alternative is to model the error covariance and transform the regression
so that its errors are homoskedastic and uncorrelated.
This leads to \gls{gls}.

\section{\glsentrylong{gls}}

Write the linear model as
\begin{align}
    Y_i=X_i\T\beta+e_i,\qquad i=1,\ldots,n,
    \qquad Y=X\beta+e,
\end{align}
where $X_i\in\reals^k$, $k<n$, $X$ has rows $X_i\T$, and $X$ has full column rank.
An intercept, if included, counts as one of the $k$ coefficients.

\begin{assumption}[Conditional Mean and Covariance]\label{ass-gls}
    Conditional on the entire design matrix $X$,
    \begin{align}
        \E(e\given X)=0,\qquad
        \var(e\given X)=\Omega,
    \end{align}
    where $\Omega$ is positive definite.
    The covariance matrix $\Omega$ may be a function of $X$.
\end{assumption}

\begin{remark}
	Under this assumption, \gls{ols} is still conditionally unbiased, but its covariance is
	\begin{align}
		\var(\widehat\beta_{\mathrm{OLS}}\given X)
		=(X\T X)\inv X\T\Omega X(X\T X)\inv.
	\end{align}
	When $\Omega$ is not proportional to $I_n$, \gls{ols} need not be efficient.
	\Gls{gls} can improve efficiency by accounting for the covariance structure.
\end{remark}

\begin{remark}[Mahalanobis Distance]
	\Gls{ols} minimizes the squared Euclidean length of the residual vector,
	$(Y-Xb)\T(Y-Xb)$.
	If some observations are much noisier than others, or their errors move together,
	this does not account for how informative each residual direction is.
	The squared \emph{Mahalanobis distance} instead scales by the inverse covariance.
\end{remark}

\begin{definition}[\glsentrylong{gls}]
    When $\Omega$ is known, the \gls{gls} estimator minimizes
    \begin{align}\label{eq-gls-objective}
        \widehat\beta_{\mathrm{GLS}}
        =\arg\min_{b\in\reals^k}(Y-Xb)\T\Omega\inv(Y-Xb).
    \end{align}
    Its first-order condition and solution are
    \begin{align}
        X\T\Omega\inv(Y-X\widehat\beta_{\mathrm{GLS}})=0
		\implies
        \widehat\beta_{\mathrm{GLS}}
        =(X\T\Omega\inv X)\inv X\T\Omega\inv Y.\label{eq-gls}
    \end{align}
\end{definition}

\begin{remark}[\glsentryshort{gls} is \glsentryshort{ols} on Transformed Data]
	Choose an invertible matrix $C$ such that $C\T C=\Omega\inv$.
	Premultiplying the \emph{entire equation} by $C$ gives
	\begin{align}
		Y^*=X^*\beta+e^*,
		\qquad Y^*=CY,\quad X^*=CX,\quad e^*=Ce.
	\end{align}
	The transformed errors satisfy
	\begin{align}
		\E(e^*\given X)=0,\qquad
		\var(e^*\given X)= C\Omega C\T= I_n.
	\end{align}
	Applying \gls{ols} to the transformed regression yields
	\begin{align}
		\bigl((X^*)\T X^*\bigr)\inv (X^*)\T Y^*
		=(X\T\Omega\inv X)\inv X\T\Omega\inv Y.
	\end{align}
	This transformation is called \emph{whitening}.
	It removes error covariances and equalizes variances;
	it does not generally make the errors independent.
	Both the outcome and every regressor must be transformed, including any intercept.
\end{remark}

\begin{theorem}[Properties of \glsentryshort{gls}]\label{thm-gls}
    Under \autoref{ass-gls} and full column rank,
    \begin{align}
        \E(\widehat\beta_{\mathrm{GLS}}\given X)&=\beta,\\
        \var(\widehat\beta_{\mathrm{GLS}}\given X)
        &=(X\T\Omega\inv X)\inv.
    \end{align}
    Moreover, \gls{gls} has the smallest conditional covariance matrix
    among estimators linear in $Y$ and conditionally unbiased for every $\beta$.
\end{theorem}
\begin{proof}[Sketch of proof]
    Substituting $Y=X\beta+e$ into \eqref{eq-gls} gives
    \begin{align*}
        \widehat\beta_{\mathrm{GLS}}-\beta
        =(X\T\Omega\inv X)\inv X\T\Omega\inv e.
    \end{align*}
    Taking the conditional mean and covariance proves the two formulas.
    The efficiency claim follows by applying the Gauss-Markov theorem to the whitened regression.
\end{proof}

\begin{remark}[What Assumptions Does Reweighting Require?]
	For consistency, \gls{ols} can work with the unconditional moment restriction
	$\E(X_ie_i)=0$ and a suitable law of large numbers.
	Reweighting changes the moment condition.
	For example, with positive weights $w(X_i)$, \gls{wls} solves
	\begin{align}
		\sum_{i=1}^n w(X_i)X_i(Y_i-X_i\T\widehat\beta_{\mathrm{WLS}})=0.
	\end{align}
	Its population target therefore depends on $\E[w(X_i)X_ie_i]$,
	which need not vanish just because $\E(X_ie_i)=0$.
	A sufficient condition is the conditional mean restriction:
	\begin{align}
		\E(e_i\given X_i)=0
		\implies
		\E[w(X_i)X_ie_i]
		=\E\{w(X_i)X_i\E(e_i\given X_i)\}=0.
	\end{align}
	If weights depend on some other variable $Z_i$, it suffices to assume $\E(e_i\given X_i,Z_i)=0$.
\end{remark}

\begin{remark}[Stronger Than Which \glsentryshort{ols} Assumption?]
    The conditional mean restriction is stronger than the unconditional orthogonality
    sufficient for \gls{ols} consistency.
    It is not stronger than the full-design conditional mean assumption already used
    to prove finite-sample \gls{ols} unbiasedness.
    Reweighting can change the population regression coefficient when the linear model
    is only a best linear approximation rather than the conditional mean.
\end{remark}

\begin{remark}[Consistency]
	More generally, consistency of \gls{gls} follows if
	\begin{align}
		\frac1n X\T\Omega\inv X\pto Q,\qquad
		\frac1n X\T\Omega\inv e\pto0,
	\end{align}
	where $Q$ is positive definite.
	The conditional mean assumption makes the weighted moments mean zero,
	but laws of large numbers and appropriate moment conditions are still needed.
\end{remark}

\subsection{Feasible \glsentryshort{gls}}

In practice, the error covariance is usually unknown.
\Gls{fgls} replaces $\Omega$ with an estimate:
\begin{align}\label{eq-fgls}
    \widehat\beta_{\mathrm{FGLS}}
    =(X\T\widehat\Omega\inv X)\inv X\T\widehat\Omega\inv Y.
\end{align}
An unrestricted $n\times n$ covariance matrix has too many entries to estimate
from one regression sample.
We therefore specify a structure, such as a variance function or an autoregressive error model.

\begin{definition}[\gls{fgls} Procedure]
	A typical \gls{fgls} procedure is
	\begin{enumerate}
		\item Obtain a consistent preliminary estimate of $\beta$, often by \gls{ols}.
		\item Use its residuals to estimate the parameters of the covariance model.
		\item Construct $\widehat\Omega$ and apply \eqref{eq-fgls}.
	\end{enumerate}
\end{definition}

\begin{remark}
	For a correctly specified, regular finite-dimensional covariance model,
	consistent estimation of its parameters and suitable moment conditions can give
	\begin{align}
		\sqrt n\bigl(\widehat\beta_{\mathrm{FGLS}}-\widehat\beta_{\mathrm{GLS}}\bigr)\pto0.
	\end{align}
	Then \gls{fgls} has the same first-order asymptotic distribution as infeasible \gls{gls}.
\end{remark}

\begin{remark}[Efficiency versus Robustness]
    Correctly modeling the covariance can improve efficiency.
    If that model is wrong, the usual \gls{gls} variance formula need not be valid,
    and the efficiency claim is lost.
\end{remark}

\section{Two Examples of \glsentryshort{gls}}

\subsection{Heteroskedasticity and Weighted Least Squares}

Suppose observations are independent, $\E(e_i\given X_i)=0$, and
$\var(e_i\given X_i)=\sigma_i^2>0$.
The covariance matrix is diagonal, so \gls{gls} becomes \gls{wls}:
\begin{align}
    \widehat\beta_{\mathrm{WLS}}
    =\arg\min_b\sum_{i=1}^n\frac{(Y_i-X_i\T b)^2}{\sigma_i^2}
    =\left(\sum_{i=1}^n\frac{X_iX_i\T}{\sigma_i^2}\right)\inv
      \sum_{i=1}^n\frac{X_iY_i}{\sigma_i^2}.
\end{align}
Equivalently, divide the entire equation by the conditional standard deviation:
\begin{align}
    \frac{Y_i}{\sigma_i}
    =\left(\frac{X_i}{\sigma_i}\right)\T\beta+\frac{e_i}{\sigma_i}.
\end{align}
The transformed error has conditional variance one.
Thus the objective weight is the inverse \emph{variance},
while the data are divided by the \emph{standard deviation}.
Holding the regressors fixed, less noisy observations receive more weight.

\begin{example}[Estimating a Common Mean]
    Let $Y_1=\mu+e_1$ and $Y_2=\mu+e_2$ have independent mean-zero errors
    with variances $1$ and $4$.
    The simple average has variance $(1+4)/4=5/4$.
    Inverse-variance weighting instead gives
    \begin{align}
        \widehat\mu_{\mathrm{WLS}}
        =\frac{Y_1+Y_2/4}{1+1/4}
        =\frac45Y_1+\frac15Y_2,
        \qquad
        \var(\widehat\mu_{\mathrm{WLS}})=\frac45.
    \end{align}
    The noisier observation still contributes, but receives less weight.
\end{example}

\begin{remark}[Estimating the Variance]
	Suppose $\E(e_i\given X_i,Z_i)=0$ and model the conditional variance as
	\begin{align}
		\E(e_i^2\given X_i,Z_i)=\exp(Z_i\T\gamma),
	\end{align}
	where $Z_i$ contains observed variance predictors.
	An intercept in $Z_i$ captures the overall scale, and the exponential ensures positive variances.
	After a preliminary consistent regression, fit this variance function to the squared residuals:
	\begin{align}
		\widehat\gamma
		=\arg\min_\gamma\frac1n\sum_{i=1}^n
		\left\{\widehat e_i^2-\exp(Z_i\T\gamma)\right\}^2.
	\end{align}
	This is simply nonlinear least squares.
	Then use $\widehat\sigma_i^2=\exp(Z_i\T\widehat\gamma)$ and weights
	$\widehat w_i=1/\widehat\sigma_i^2$ in the weighted regression.
	Consistency requires a correctly specified and identified variance function,
	suitable fourth moments, and the usual regularity conditions.
\end{remark}

\subsection{AR(1) Errors and Quasi-Differencing}

Consider the time series regression
\begin{align}\label{eq-ar-errors}
    Y_t=X_t\T\beta+e_t,\qquad
    e_t=\rho e_{t-1}+v_t,\qquad |\rho|<1.
\end{align}
For a simple sufficient set of assumptions, let $\{v_t\}$ be \gls{iid}
with mean zero and variance $\sigma_v^2>0$,
independent of the entire regressor process.
Use the stationary causal solution for $e_t$.
Then
\begin{align}
    \Omega_{ts}=\cov(e_t,e_s\given X)
    =\frac{\sigma_v^2}{1-\rho^2}\rho^{|t-s|}.
\end{align}
The errors are correlated across observations, so diagonal weights alone cannot whiten them.

\begin{remark}[\glsentryshort{ols} Is Consistent but Generally Inefficient]
    Also assume that $\{(X_t,e_t)\}$ is stationary and ergodic,
    $\E\|X_t\|^2<\infty$, and $Q_{xx}=\E(X_tX_t\T)$ is positive definite.
    Independence of the regressor and innovation processes implies $\E(X_te_t)=0$.
    The law of large numbers then gives
    \begin{align}
        \widehat\beta_{\mathrm{OLS}}-\beta
        &=\left(\frac1T\sum_{t=1}^T X_tX_t\T\right)\inv
          \left(\frac1T\sum_{t=1}^T X_te_t\right)
        \pto Q_{xx}\inv\E(X_te_t)=0.
    \end{align}
    Thus serial correlation in the errors does not itself make \gls{ols} inconsistent.
    However, \gls{ols} does not generally use the covariance structure efficiently.
    By \autoref{thm-gls}, \gls{gls} has a weakly smaller conditional covariance matrix.
\end{remark}

\subsubsection{Transforming the Regression}

For known $\rho$, subtract $\rho$ times the equation at $t-1$ from the equation at $t$:
\begin{align}\label{eq-quasidifference}
    \underbrace{Y_t-\rho Y_{t-1}}_{Y_t^*}
    =\underbrace{(X_t-\rho X_{t-1})\T}_{(X_t^*)\T}\beta+v_t,
    \qquad t=2,\ldots,T.
\end{align}
This is \emph{quasi-differencing}.
It removes the serial correlation in the errors, and \gls{ols} on the transformed equation
implements \gls{gls} apart from the treatment of the first observation.
For exact \gls{gls} with stationary errors, also use
\begin{align}
    Y_1^*=\sqrt{1-\rho^2}\,Y_1,\qquad
    X_1^*=\sqrt{1-\rho^2}\,X_1.
\end{align}
Its error is $\sqrt{1-\rho^2}\,e_1$, which has variance $\sigma_v^2$
and is uncorrelated with $v_2,\ldots,v_T$.
The transformed error covariance is therefore $\sigma_v^2 I_T$.
Dividing all transformed observations by $\sigma_v$ would give unit error variances,
but this common rescaling does not change the coefficient estimate.
If $\rho$ is fixed away from the unit-root boundary,
dropping this first observation is asymptotically negligible under the usual regularity conditions.

\begin{remark}[Transform the Intercept Too]
    If $X_t$ contains a constant one, the corresponding transformed regressor is $1-\rho$
    for $t\ge2$ and $\sqrt{1-\rho^2}$ for the first observation.
    Transforming only $Y_t$, or leaving the original regressor columns unchanged,
    does not implement \gls{gls}.
\end{remark}

\begin{remark}[Why Exogeneity Matters]
	For consistency of the transformed regression, we need a law of large numbers,
	a nonsingular limiting regressor second moment, and
	\begin{align}\label{eq-transformed-moment}
		\E\bigl[(X_t-\rho X_{t-1})v_t\bigr]=0.
	\end{align}
	Contemporaneous orthogonality $\E(X_te_t)=0$ alone does not imply this.
	Expanding the transformed moment gives
	\begin{align}
		\E\bigl[(X_t-\rho X_{t-1})(e_t-\rho e_{t-1})\bigr]
		={}&\E(X_te_t)-\rho\E(X_te_{t-1})\\
		&-\rho\E(X_{t-1}e_t)+\rho^2\E(X_{t-1}e_{t-1}).
	\end{align}
	The cross-date terms need not vanish.
	A sufficient condition is \emph{strict exogeneity},
	$\E(e_t\given\{X_s\}_{s\in\mathbb Z})=0$ for every $t$.
	Our independence assumption above implies this condition.
	It ensures that the error remains mean zero conditional on the full transformed regressor history.
\end{remark}

\begin{remark}[Sufficient, Not Necessary]
    Strict exogeneity is a convenient sufficient condition, not a necessary one.
    Some dynamic models satisfy \eqref{eq-transformed-moment} without it.
    In such cases, consistency must be checked directly using the transformed regressors and errors;
    whitening the errors by itself is not a proof.
\end{remark}

\begin{remark}[\gls{fgls} for AR(1) Errors]
	When the original regression admits a consistent \gls{ols} estimate, form its residuals
	and estimate the error autocorrelation by
	\begin{align}
		\widehat\rho
		=\frac{\sum_{t=2}^T\widehat e_t\widehat e_{t-1}}
		{\sum_{t=2}^T\widehat e_{t-1}^2}.
	\end{align}
	Under stationarity, ergodicity, and suitable moments, $\widehat\rho\pto\rho$.
	Replace $\rho$ with $\widehat\rho$ in the transformation and estimate $\beta$ by \gls{ols}.
	For a stationary fit, constrain the covariance estimate to $|\widehat\rho|<1$.
	Using all $T$ transformed observations, estimate the innovation variance by
	\begin{align}
		\widehat\sigma_v^2
		=\frac{1}{T-k}\sum_{t=1}^T
		\bigl(Y_t^*-(X_t^*)\T\widehat\beta_{\mathrm{FGLS}}\bigr)^2.
	\end{align}
	Under correct specification and the usual feasible-\gls{gls} conditions,
	the model-based covariance estimator is
	\begin{align}
		\widehat\sigma_v^2
		\left(\sum_{t=1}^T X_t^*(X_t^*)\T\right)\inv.
	\end{align}
	If the preliminary regression is inconsistent because of endogeneity,
	its residuals do not generally yield a consistent estimate of $\rho$.
	The procedure is therefore not a general cure for endogenous regressors.
\end{remark}

\vfill
\section*{Acronyms}
\printglossaries

\end{document}
